\documentclass[11pt]{article}
\usepackage[margin=1in]{geometry}
\usepackage{amsmath,amssymb,amsthm}
\usepackage{graphicx}
\usepackage[hidelinks]{hyperref}

\newtheorem{theorem}{Theorem}

\title{Affirmative Answer for the BDP Status of the Operator \(S\) in arXiv:2209.01981}
\author{Run \texttt{fa\_banach\_001}, lane 13}
\date{2026-06-16}

\begin{document}
\maketitle

\section*{Status}

Candidate full solution, likely valid. This packet answers the local question
asked after Example 3.1.1 of Emelyanov--Gorokhova \cite{EG2022}: whether the
operator \(S:L^1[0,1]\to \ell^\infty\) defined there is a \([BDP]\)-operator.

\section*{Source Question}

The source paper defines
\[
  S(f)=\left(\int_0^1 f(t)r_k^+(t)\,dt\right)_{k=1}^{\infty},
  \qquad f\in L^1[0,1],
\]
where \(r_k\) are the Rademacher functions. After Example 3.1.1, the authors
write:
\[
  \text{``We do not know whether or not the operator }S
  \text{ in Example 3.1.1 is a [BDP]-operator.''}
\]
The relevant source crop from page 15 is shown below.

\begin{center}
\includegraphics[width=\textwidth]{figures/open_problem_crop.png}
\end{center}

\section*{Definitions Used}

Following Definition 1.2.2 of \cite{EG2022}, an operator \(T:X\to Y\) is a
\([GPP]\)-operator if \(T\) carries limited subsets of \(X\) to relatively
compact subsets of \(Y\). It is a \([BDP]\)-operator if \(T\) carries limited
subsets of \(X\) to relatively weakly compact subsets of \(Y\).

The source paper also recalls the classical fact of Bourgain--Diestel
\cite{BD1984} that all separable Banach spaces have the Gelfand--Phillips
property: every limited subset of such a space is relatively compact.

\section*{Proof Intuition}

The question is resolved before one uses any special structure of \(S\). The
domain \(L^1[0,1]\) is separable. Therefore its limited subsets are already
relatively compact in \(L^1[0,1]\). A bounded linear operator sends relatively
compact sets to relatively compact sets, and relative compactness is stronger
than relative weak compactness. Hence every bounded operator out of
\(L^1[0,1]\), including the displayed operator \(S\), is automatically
\([GPP]\) and therefore \([BDP]\).

\begin{theorem}
The operator
\[
  S:L^1[0,1]\to \ell^\infty,\qquad
  S(f)=\left(\int_0^1 f(t)r_k^+(t)\,dt\right)_{k=1}^{\infty},
\]
from Example 3.1.1 of \cite{EG2022} is a \([BDP]\)-operator. More strongly,
it is a \([GPP]\)-operator.
\end{theorem}

\begin{proof}
The space \(L^1[0,1]\) is separable. By the Bourgain--Diestel theorem recalled
in \cite{EG2022}, every separable Banach space has the Gelfand--Phillips
property. Thus every limited subset \(A\subseteq L^1[0,1]\) is relatively
compact in \(L^1[0,1]\).

Let \(T:L^1[0,1]\to Y\) be any bounded linear operator into a Banach space
\(Y\), and let \(A\subseteq L^1[0,1]\) be limited. Since \(A\) is relatively
compact and \(T\) is continuous, \(T(A)\) is relatively compact in \(Y\).
Therefore \(T\) is a \([GPP]\)-operator.

Every relatively compact subset of a Banach space is relatively weakly
compact, so every \([GPP]\)-operator is a \([BDP]\)-operator. Applying this
to \(T=S\) gives the desired conclusion.
\end{proof}

\section*{Verification and Scope}

The argument proves a slightly stronger statement: every bounded linear
operator from \(L^1[0,1]\) to any Banach space is \([GPP]\), hence \([BDP]\).
No positivity, domination, or Rademacher-specific computation is required.

This also explains the apparent tension in Example 3.1.1 of the source paper.
The Rademacher sequence used there is described as an almost limited subset of
\(L^1[0,1]\). Such a sequence can obstruct the strong operator class
\([s\text{-}GPP]\), which tests almost limited subsets, but it is not by
itself an obstruction to \([GPP]\), whose definition tests limited subsets.
Indeed, since \(L^1[0,1]\) is separable and has the Gelfand--Phillips property,
any limited subset of \(L^1[0,1]\) must be relatively compact.

\section*{Novelty Check}

Before promotion, the run indexes were searched for arXiv id 2209.01981 and
the phrases ``Operators Affiliated'', ``Banach lattice properties'', and
``BDP''. Web searches were also made for the exact and near-exact phrases
\[
\texttt{"We do not know whether or not the operator S" "BDP"},
\]
\[
\texttt{"Operators Affiliated to Banach Lattice Properties" "BDP" "S"},
\]
and
\[
\texttt{"[BDP]-operator" "r\_k\^+"}.
\]
No separate later paper answering this exact question was found. The result is
therefore recorded as a direct solution extracted from the source definitions
and the standard Bourgain--Diestel fact recalled in the source paper.

\section*{Human Review Recommendation}

Reviewers should check only that the source definitions of \([GPP]\) and
\([BDP]\) are being used exactly as printed, namely for limited subsets of
the domain. Under those definitions, the proof is immediate. If the authors
intended the strong almost-limited class instead, then the source question
should be relabelled as a question about \([s\text{-}BDP]\), not \([BDP]\).

\begin{thebibliography}{9}

\bibitem{EG2022}
E. Emelyanov and S. Gorokhova,
\emph{Operators Affiliated to Banach Lattice Properties and the Enveloping Norms},
arXiv:2209.01981, 2022.
\url{https://arxiv.org/abs/2209.01981}

\bibitem{BD1984}
J. Bourgain and J. Diestel,
\emph{Limited operators and strict cosingularity},
Math. Nachr. 119 (1984), 55--58.

\end{thebibliography}

\end{document}

